\scanpage{21}
\hfill\fbox{Motifs} % L1

1) La catégorie $M^+(X)$ % L2

2) Variance \unsure{avec|pour} $X$ % L3

3) Cas $X=\varprojlim X_i$ % L4

4) Foncteurs $T_\ell$ % L5

5) Les $\mathbb Q_\ell(-n)$ % L6

6) La catégorie $M(X)$ % L7

7) Les foncteurs $\mathrm{Hom}$ et $\mathrm{RHom}$ % L8

8) Motifs constants, tordus, et \unsure{pl. con.|pt. con.|pot. con.}. % L9

9) Filtrations de $M(X)$ et $M^+(X)$ % L10

10) Anneaux $M^+(X)$ et $M(X)$. % L11

11) Interprétation topologique des \deleted{\unsure{types de niveau}} \inserted{niveaux} % L12
par la filtration canonique \deleted{\unsure{de leurs}} variétés proj. non singulièr\unsure{es|e}. % L13

12) L'hom $L(K)\to M^+(K)$ et invariants birationnels fondamentaux. % L14
[N.B. À généraliser à des coefficients quelconques] % L15

13) Caractérisation galoisienne des filtrations (conj. de Tate) % L16

14) Invariants de Galois et \unsure{thm.|thms} de commutation. % L17

15) Cohomologie absolue des motifs. Cas des schémas de type fini. % L18

16) Relation avec les pts rationnels et la coh. des VA sur % L19
un schéma de type fini. % L20

17) Formes positives. % L21

18) Motifs et fonctions $L$ % L22

19) Relations avec la cohomologie de Hodge ....  % L23

% archive sheet number: 20
